\section{Reacciones}

\begin{align*}
    \text{Reacción reversible}
    \\
    aA + bB \rightarrow cC + dD
    \\
    \\
    \text{Reacción irreversible}
    \\
    aA + bB ⇌ cC + dD
\end{align*}

\section{Constante de concentración ($K$)}

\begin{equation*}
    K=\frac {[C]^c*[D]^d}{[A]^a*[B]^b}
\end{equation*}

Los sólidos y líquidos no cambian en su [x] inicial y final, por lo tanto sus valores son 1 y por ello no se escriben en la ecuación de equilibrio.

\section{Ácidos y bases según Bronsten-Lowry}

\begin{itemize}
    \item Ácidos: Donan protones $[H^{+}]$
    \item Bases: Aceptan protones $[H^{+}]$ 
\end{itemize}

Se denominan ácidos fuertes a quellos ácidos que se desprotonan de manera completa. Estos son: 

\begin{itemize}
    \item $HCl$
    \item $HNO_{3}$
    \item $H_{2}SO_{4}$
    \item $HBr$
    \item $HI$
    \item $HClO:{4}$
\end{itemize}

El resto de ácidos se denominan ácidos débiles

\section{Valores de Kw y pH}

\begin{align*}
    Kw &= [H_{3}O^{+}]*[OH^{-}]
    \\
    Kw &= 1*10^{-14}
    \\
    pH &= -log[H_{3}O^{+}]
\end{align*}
